% Ebenen – bank/ebenen/ (zusammenbau v0.4, Heft)
\einheitenkopf[t4]{Ebenen}
\setcounter{aufgabe}{36}

% Anlauf (ohne Prüfkennung)
\begin{aufgabe}{Parameterform – Anlauf}
\begin{teile}
\teil Eine Ebene soll durch $A(1 | 3 | 0)$, $B(2 | 0 | 4)$ und $C(5 | 1 | 1)$ gehen. Was ist gegeben? \\ \kreuz{drei Punkte} \\ \kreuz{ein Punkt und zwei Richtungen} \\ \kreuz{eine Gerade und ein Punkt} \\ \kreuz{zwei Geraden}
\teil Punkt $(2 | 1 | 3)$, Spannvektoren $(1 | 0 | 4)$ und $(0 | 3 | 1)$: Parametergleichung der Ebene $E$?
\teil $A(2 | 1 | -1)$, $B(4 | 2 | 0)$, $C(0 | 3 | 1)$: Parametergleichung der Ebene durch $A$, $B$, $C$? Prüfe, dass die Spannvektoren nicht kollinear sind.
\end{teile}
\end{aufgabe}

\begin{aufgabe}{Parameterform – Prüfungsaufgaben}
\begin{teile}
\teil Ein Zeltdach hat vier rautenförmige Dachflächen; eine davon hat die Ecken $C(6 | 6 | 0)$, $F(6 | 3 | 3)$, $G(3 | 6 | 3)$ und die Spitze $S$ (1 LE = 1 m). Gib eine Parametergleichung der Ebene $L$ durch $C$, $F$ und $G$ an und weise nach, dass $S(3 | 3 | 6)$ in $L$ liegt. \hfill (Abitur 2022 GK)
\teil Gegeben sind $A(2 | -1 | 1)$, $B(1 | 3 | 4)$ und der Punkt $C$ mit $\overrightarrow{OC} = 3 \cdot \overrightarrow{OA}$; $O$ ist der Koordinatenursprung. Begründe, dass es genau eine Ebene gibt, die $O$, $A$, $B$ und $C$ enthält. \hfill (Abitur 2018 GK)
\teil Gegeben sind $A(-1 | 2 | 2)$, $B(4 | 0 | 1)$ und der Punkt $C$ mit $\overrightarrow{OC} = -2 \cdot \overrightarrow{OA}$; $O$ ist der Koordinatenursprung. Begründe, dass es genau eine Ebene gibt, die $O$, $A$, $B$ und $C$ enthält. \hfill (Abitur 2018 GK)
\teil Das Pultdach eines Carports ist das Rechteck mit den Ecken $E(1 | 0 | 3)$, $F(7 | 0 | 3)$, $G(7 | 5 | 6)$ und $H(1 | 5 | 6)$ (1 LE = 1 m). Gib eine Parametergleichung der Dachebene $L$ an und prüfe, ob der Befestigungspunkt $P(4 | 2{,}5 | 4{,}5)$ einer Lampe in $L$ liegt. \hfill (Abitur 2022 GK)
\teil Die Dachfläche eines Gartenhauses ist das Rechteck mit den Ecken $A(0 | 0 | 3)$, $B(6 | 0 | 3)$, $C(6 | 4 | 5)$ und $D(0 | 4 | 5)$ (1 LE = 1 m). Gib eine Parametergleichung der Dachebene $L$ an und prüfe, ob der Punkt $P(3 | 3 | 4)$ in $L$ liegt. \hfill (Abitur 2022 GK)
\end{teile}
\end{aufgabe}

% Anlauf (ohne Prüfkennung)
\begin{aufgabe}{Koordinatengleichung bestimmen – Anlauf}
\begin{teile}
\teil Gesucht ist eine Gleichung der Form $ax + by + cz = d$ für die Ebene durch drei gegebene Punkte. Welche Form ist verlangt? \\ \kreuz{Parameterform} \\ \kreuz{Koordinatenform} \\ \kreuz{Normalenform}
\teil Spannvektoren $(1 | 0 | 3)$ und $(0 | 1 | 2)$: ein Normalenvektor $\vec n$ der Ebene? $\vec n = ($ \leerfeld | \leerfeld | \leerfeld $)$
\teil Eine Ebene schneidet die Achsen in $(6 | 0 | 0)$, $(0 | 3 | 0)$ und $(0 | 0 | 2)$. Koordinatengleichung?
\end{teile}
\end{aufgabe}

\begin{aufgabe}{Koordinatengleichung bestimmen – Prüfungsaufgaben}
\begin{teile}
\teil Ein dreieckiges Sonnensegel ist an drei Masten in $A(3 | 2 | 2)$, $B(3 | 5 | 2)$ und $C(0 | 3 | 4)$ befestigt (1 LE = 1 m). Die Ebene des Segels hat eine Gleichung der Form $2x + 3z = d$. Bestimme $d$. \hfill (Abitur 2020 GK) \\ $d = $ \leerfeld
\teil Ein Dachaufbau hat die Form einer Pyramide mit der Grundfläche $A(1 | 1 | 0)$, $B(4 | 1 | 0)$, $C(4 | 4 | 0)$, $D(1 | 7 | 0)$ und der Spitze $S(1 | 1 | 5)$ (1 LE = 1 m). Die Seitenfläche $CDS$ liegt in der Ebene $E$. Bestimme eine Gleichung von $E$ in Koordinatenform (zur Kontrolle: $5x + 5y + 6z = 40$). \hfill (Abitur 2024 GK)
\teil Die Seitenfläche $BCS$ eines pyramidenförmigen Kristalls hat die Ecken $B(3 | 1 | 3)$, $C(2 | 4 | 3)$ und $S(1 | 1 | 6)$ (1 LE = 1 mm). Bestimme eine Koordinatengleichung der Ebene $E$, in der $BCS$ liegt (zur Kontrolle: $3x + y + 2z = 16$). \hfill (Abitur 2025 GK)
\teil Ein Holzkörper ist eine Pyramide über dem Rechteck $ABCD$ mit $B(6 | 0 | 0)$ und $C(6 | 8 | 0)$; seine Spitze $E(0 | 8 | 4)$ liegt senkrecht über $D$ (1 LE = 1 cm). Bestimme eine Koordinatengleichung der Ebene $L$ durch $B$, $C$ und $E$. \hfill (Abitur 2021 GK)
\teil Bestimme eine Koordinatengleichung von $E\colon \vec x = (4 | 1 | 0) + r \cdot (-2 | 4 | 0) + s \cdot (-2 | 0 | 3)$, $r, s \in \mathbb{R}$ (zur Kontrolle: $6x + 3y + 4z = 27$). \hfill (Abitur 2020 GK)
\teil Ein Sonnensegel $W$ ist in $M(3 | 3 | 2)$, $F(0 | 0 | 3)$ und $S(6 | 0 | 0)$ befestigt (1 LE = 1 m). Bestimme eine Koordinatengleichung von $W$ mit dem Ansatz $ax + by + cz = 18$ (zur Kontrolle: $3x - y + 6z = 18$). \hfill (Abitur 2024 GK)
\teil Eine dreieckige Plane über einer Lagerhalle hat die Ecken $U(1 | 1 | 5)$, $V(3 | 9 | 5)$ und $W(9 | 1 | 3)$ (1 LE = 1 m). Gib die Ebene $E$ der Plane in Parameter- und in Koordinatenform an (zur Kontrolle: $4x - y + 16z = 83$). \hfill (Abitur 2020 GK)
\teil Die seitlichen Streben einer Brücke liegen in der Ebene $E$ durch $A(1 | 0 | 6)$, $B(3 | 40 | 6)$ und $C(1{,}2 | 4 | 8)$; die Fahrbahn liegt in der $xy$-Ebene (1 LE = 1 m). Bestimme einen Normalenvektor mit dem Vektorprodukt und gib $E$ in Parameter- und in Koordinatenform an. \hfill (Abitur 2018 GK)
\end{teile}
\end{aufgabe}

% Anlauf (ohne Prüfkennung)
\begin{aufgabe}{Normalenvektor und Nachweise – Anlauf}
\begin{teile}
\teil Gegeben ist $E\colon \vec x = \vec p + r \cdot \vec u + s \cdot \vec v$ mit Zahlen für die Vektoren. In welcher Form ist $E$ gegeben? \\ \kreuz{Parameterform} \\ \kreuz{Koordinatenform} \\ \kreuz{Normalenform}
\teil Gib zwei Normalenvektoren von $E\colon 4x - 2y + z = 7$ an.
\teil Gib $E\colon \left(\vec x - (1 | -2 | 0)\right) \cdot (2 | 1 | 5) = 0$ in Koordinatenform an.
\end{teile}
\end{aufgabe}

\begin{aufgabe}{Normalenvektor und Nachweise – Prüfungsaufgaben}
\begin{teile}
\teil Gegeben ist $F\colon \left(\vec x - (2 | 5 | 1)\right) \cdot (-3 | 1 | 4) = 0$. Gib $F$ in Koordinatenform an. \hfill (Abitur 2021 GK)
\teil Gegeben sind $E(5 | 2 | 3)$, $F(1 | 4 | 3)$ und $S(1 | 1 | 5)$. Bestimme rechnerisch einen Normalenvektor der Ebene $EFS$ und prüfe ihn. \hfill (Abitur 2025 GK)
\end{teile}
\end{aufgabe}

% Anlauf (ohne Prüfkennung)
\begin{aufgabe}{Besondere Lagen – Anlauf}
\begin{teile}
\teil $E\colon 4x + 5z = 30$. Welche Variable fehlt? \\ \kreuz{$x$} \\ \kreuz{$y$} \\ \kreuz{$z$} \\ \kreuz{keine}
\teil Gleichung der $xz$-Ebene?
\teil Zu welcher Achse ist $E\colon 5x + z = 5$ parallel? Begründe.
\end{teile}
\end{aufgabe}

\begin{aufgabe}{Besondere Lagen – Prüfungsaufgaben}
\begin{teile}
\teil Welche besondere Lage hat $E\colon 4x - 5y = 0$? Beschreibe sie. \hfill (Abitur 2022 GK)
\teil Eine Kletterwand ist ein ebenes Viereck mit $A(2 | 6 | 1)$, $B(6 | 4 | 1)$, $C(5 | 4{,}5 | 5)$ und $D(1 | 6{,}5 | 5)$; die $xy$-Ebene ist der Boden (1 LE = 1 m). Weise nach, dass die Wand vertikal steht. \hfill (Abitur 2022 GK)
\end{teile}
\end{aufgabe}

% Anlauf (ohne Prüfkennung)
\begin{aufgabe}{Parallele Ebenen – Anlauf}
\begin{teile}
\teil $E\colon x + 2y - z = 3$ und $F\colon 2x + 4y - 2z = 6$. Wie liegen sie zueinander? \\ \kreuz{parallel} \\ \kreuz{identisch} \\ \kreuz{schneidend}
\teil Ebene $F$ parallel zu $E\colon 3x + y - 2z = 4$ durch $P(2 | 1 | 1)$?
\teil $F$ ist parallel zu $E\colon 2x - y + 2z = 4$ und geht durch $P(1 | 1 | 3)$. Liegt $Q(3 | 1 | 1)$ in $F$?
\end{teile}
\end{aufgabe}

\begin{aufgabe}{Parallele Ebenen – Prüfungsaufgaben}
\begin{teile}
\teil Ein Würfel mit der Kantenlänge 6 hat eine Ecke im Ursprung und liegt im ersten Oktanten. Die Punkte $K(0 | 6 | 3)$ und $L(1 | 0 | 6)$ liegen auf Kanten. $K$ liegt in einer Ebene $T$, die parallel zu $S\colon 3x + 2y + 3z = 0$ ist. Liegt auch $L$ in $T$? \hfill (Abitur 2019 GK)
\teil Bei einem Körper liegen die Ecken $A$, $B$, $C$ in der Ebene $L_1\colon 2x + y + 2z = 8$, die Ecken $D$, $E$, $F$ in $L_2\colon 4x + 2y + 4z = 14$. Begründe, dass $L_1$ und $L_2$ keine gemeinsamen Punkte haben. \hfill (Abitur 2023 GK)
\teil Eine Halle hat die quadratische Grundfläche $ABCD$ mit 6 m Seitenlänge in der $xy$-Ebene, $A$ im Ursprung, $B$ auf der $x$-Achse, $D$ auf der $y$-Achse (1 LE = 1 m). Ein Zuschauerboden liegt in der Ebene $E_B$, die parallel zu $P\colon x - y + 6z = 30$ durch $B$ verläuft. Gib $E_B$ an und bestimme die größte Höhe des Bodens über der Grundfläche. \hfill (Abitur 2020 GK)
\teil Ein dreiseitiges Prisma hat die Grundfläche $ABC$ in der Ebene $L\colon x + 2y + 2z = 3$ und die Deckfläche $DEF$; die Seitenkante $BE$ verbindet $B(1 | 1 | 0)$ mit $E(4 | 4 | 6)$. Eine zu $L$ parallele Ebene $M$ teilt das Prisma so, dass der Teilkörper, der $E$ enthält, doppelt so groß ist wie der andere. Bestimme eine Gleichung von $M$. \hfill (Abitur 2020 GK)
\teil Ein Prisma hat die Grundfläche in der Ebene $L\colon 2x - y + 2z = 1$; eine Seitenkante verbindet $A(1 | 1 | 0)$ in der Grundfläche mit $D(5 | 5 | 4)$ in der Deckfläche. Eine zu $L$ parallele Ebene $M$ teilt das Prisma so, dass der Teilkörper, der $D$ enthält, dreimal so groß ist wie der andere. Bestimme eine Gleichung von $M$. \hfill (Abitur 2020 GK)
\end{teile}
\end{aufgabe}

% Anlauf (ohne Prüfkennung)
\begin{aufgabe}{Parallele Ebenen über die Normalenvektoren auswählen und über die Konstante der Koordinatengleichung der Höhe nach zuordnen – Anlauf}
\begin{teile}
\teil Schräge Regalböden liegen in Ebenen ($z$ zeigt nach oben): $E_1\colon x + 3z = 6$, $E_2\colon x + 3z = 9$, $E_3\colon 2x + 6z = 30$, $E_4\colon x + 2z = 9$. Welche sind parallel zu $E_1$, und in welcher Reihenfolge liegen sie von unten nach oben?
\end{teile}
\end{aufgabe}

\begin{aufgabe}{Parallele Ebenen über die Normalenvektoren auswählen und über die Konstante der Koordinatengleichung der Höhe nach zuordnen – Prüfungsaufgaben}
\begin{teile}
\teil Ein Turm hat drei Dachetagen; die Dachflächen der Südseite sind parallel ($z$ zeigt nach oben, 1 LE = 1 m). Die untere liegt in $E\colon 2x + 5z = 40$, die mittlere und die obere in zwei der Ebenen I: $2x + 4z = 40$, II: $2x + 5z = 30$, III: $2x + 5z = 50$, IV: $4x + 5z = 40$, V: $2x + 5z = 60$. Ordne zu und begründe. \hfill (Abitur 2017 GK)
\end{teile}
\end{aufgabe}
